% ============================================================
%  TERM PROJECT – DOCUMENTATION
%  Databases · THGA Bochum
%
%  Blutspende-Verwaltungssystem (Blood Donation Management System)
% ============================================================
\documentclass[11pt,a4paper]{article}

\usepackage{thga-db}

\renewcommand{\thgaKurs}{Databases}
\renewcommand{\thgaDatum}{Summer Term 2026}

\begin{document}
\thispagestyle{empty}
\thgaTitel{Term Project -- Documentation}%
{Blutspende-Verwaltungssystem: a blood donation management system}

\begin{infobox}
\textbf{Author:} Amen Dabbech \\
\textbf{Module:} Introduction to Database Management Systems · THGA Bochum \\
\textbf{Stack:} PostgreSQL · FastAPI · Docker Compose · Python CLI (\texttt{.deb}) \\
\textbf{Repository:} \url{https://github.com/amendabbech/DBMS_10}
\end{infobox}

\tableofcontents
\vspace{1em}

% ============================================================
\section{Introduction and motivation}
% ============================================================

Blood donation centers often coordinate donors and appointments manually or
with simple spreadsheets. This makes it hard to enforce basic safety rules,
such as the mandatory waiting period between two donations by the same donor,
and to answer questions such as \emph{``which donors are currently eligible to
donate again?''} without manual checking.

\textbf{Blutspende-Verwaltungssystem} replaces this manual process with a
proper relational system. It manages \textbf{Spender} (donors),
\textbf{Spendezentren} (donation centers) and \textbf{Spenden} (donations),
exposes the data through a REST API, and offers a small command-line frontend
for staff.

\begin{beispielbox}
\textbf{Usage scenario.} A donor arrives at a donation center to give blood.
Staff open the client, look up the donor, and register a new donation with
date, blood type and amount. The system automatically rejects the donation if
the donor's last donation was less than 56 days ago, enforcing the safety rule
without any manual date calculation.
\end{beispielbox}

% ============================================================
\section{Data model}
% ============================================================

The domain has three entity types. A \textbf{Spender} can give many
\textbf{Spenden} over time at different \textbf{Spendezentren}, and a
Spendezentrum can host donations from many donors -- an \texttt{N:M}
relationship between Spender and Spendezentrum that is resolved by the
\textbf{Spende} entity.

\begin{beispielbox}
\textbf{ER model (textual).}
\begin{itemize}[leftmargin=1.4em, topsep=2pt]
  \item \textbf{Spender} (1) \;--\; gives \;--\; (N) \textbf{Spende}
  \item \textbf{Spendezentrum} (1) \;--\; hosts \;--\; (N) \textbf{Spende}
\end{itemize}
Thus \textbf{Spender} \texttt{N:M} \textbf{Spendezentrum}, resolved through
\textbf{Spende}.
\end{beispielbox}

\subsection{Relational schema}

Derived from the ER model, with primary keys underlined and foreign keys in
italics:

\begin{itemize}[leftmargin=1.4em]
  \item \texttt{spender}(\underline{spender\_id}, vorname, nachname,
        geburtsdatum, blutgruppe, email)
  \item \texttt{spendezentrum}(\underline{zentrum\_id}, name, adresse, stadt)
  \item \texttt{spende}(\underline{spende\_id}, \emph{spender\_id},
        \emph{zentrum\_id}, spende\_datum, menge\_ml)
\end{itemize}

\textbf{Normalisation.} Every non-key attribute depends only on the whole
primary key of its table, and there are no transitive dependencies: a donor's
email lives only in \texttt{spender}, a center's address only in
\texttt{spendezentrum}, and a donation references donors and centers by their
keys rather than duplicating their data. The schema is therefore in
\textbf{3NF}.

% ============================================================
\section{Database (PostgreSQL)}
% ============================================================

The schema is created by an init script that PostgreSQL runs on first start,
via a Docker volume mount. In addition to the three tables, a
\textbf{trigger} enforces the core business rule of the domain: a donor may
not donate again within 56 days of their last donation.

\begin{lstlisting}[language=SQL, caption={schema.sql (excerpt) -- tables and 56-day trigger}]
CREATE TABLE spender (
    spender_id SERIAL PRIMARY KEY,
    vorname VARCHAR(100) NOT NULL,
    nachname VARCHAR(100) NOT NULL,
    geburtsdatum DATE NOT NULL,
    blutgruppe VARCHAR(3) NOT NULL,
    email VARCHAR(255) UNIQUE NOT NULL
);

CREATE TABLE spendezentrum (
    zentrum_id SERIAL PRIMARY KEY,
    name VARCHAR(150) NOT NULL,
    adresse VARCHAR(255) NOT NULL,
    stadt VARCHAR(100) NOT NULL
);

CREATE TABLE spende (
    spende_id SERIAL PRIMARY KEY,
    spender_id INTEGER NOT NULL REFERENCES spender(spender_id),
    zentrum_id INTEGER NOT NULL REFERENCES spendezentrum(zentrum_id),
    spende_datum DATE NOT NULL,
    menge_ml INTEGER NOT NULL CHECK (menge_ml > 0)
);

CREATE OR REPLACE FUNCTION pruefe_spendeabstand()
RETURNS TRIGGER AS $$
BEGIN
    IF EXISTS (
        SELECT 1 FROM spende
        WHERE spender_id = NEW.spender_id
        AND spende_datum > NEW.spende_datum - INTERVAL '56 days'
        AND spende_datum < NEW.spende_datum + INTERVAL '56 days'
    ) THEN
        RAISE EXCEPTION 'Mindestabstand von 56 Tagen nicht eingehalten';
    END IF;
    RETURN NEW;
END;
$$ LANGUAGE plpgsql;

CREATE TRIGGER trg_spendeabstand
BEFORE INSERT ON spende
FOR EACH ROW
EXECUTE FUNCTION pruefe_spendeabstand();
\end{lstlisting}

% ============================================================
\section{REST API (FastAPI)}
% ============================================================

The API is the only component that talks to the database. Read endpoints are
public; write endpoints require the header \texttt{X-API-Key}.

\renewcommand{\arraystretch}{1.2}
\begin{tabularx}{\linewidth}{@{}p{1.6cm}p{4.4cm}Xc@{}}
\toprule
\textbf{Method} & \textbf{Path} & \textbf{Purpose} & \textbf{Key} \\
\midrule
\texttt{GET}  & \texttt{/spender} & list all donors & -- \\
\texttt{POST} & \texttt{/spender} & register a new donor & yes \\
\texttt{GET}  & \texttt{/spendezentrum} & list all donation centers & -- \\
\texttt{POST} & \texttt{/spendezentrum} & register a new center & yes \\
\texttt{GET}  & \texttt{/spende} & list all donations & -- \\
\texttt{POST} & \texttt{/spende} & register a new donation & yes \\
\bottomrule
\end{tabularx}

\medskip
The \texttt{POST /spende} endpoint demonstrates how a database-level
\textbf{constraint} (the 56-day trigger) is surfaced as a clean API error
instead of a raw database exception:

\begin{lstlisting}[language=Python, caption={POST /spende -- catching the trigger's exception}]
@app.post("/spende")
def create_spende(spende: SpendeIn, auth: None = Depends(check_api_key)):
    conn = get_conn()
    cur = conn.cursor(cursor_factory=RealDictCursor)
    try:
        cur.execute(
            "INSERT INTO spende (spender_id, zentrum_id, spende_datum, menge_ml) "
            "VALUES (%s, %s, %s, %s) RETURNING *",
            (spende.spender_id, spende.zentrum_id,
             spende.spende_datum, spende.menge_ml)
        )
        new_row = cur.fetchone()
        conn.commit()
        return new_row
    except psycopg2.Error as e:
        conn.rollback()
        raise HTTPException(status_code=400, detail=str(e))
    finally:
        cur.close()
        conn.close()
\end{lstlisting}

The write path is guarded by a simple dependency that checks the API key:

\begin{lstlisting}[language=Python, caption={API-key guard}]
API_KEY = "geheimschluessel123"

def check_api_key(x_api_key: str = Header(...)):
    if x_api_key != API_KEY:
        raise HTTPException(status_code=401, detail="Ungueltiger API-Key")
\end{lstlisting}

% ============================================================
\section{Frontend (Python CLI, packaged as .deb)}
% ============================================================

The frontend is a small command-line client with a numbered menu offering six
actions: list and create Spender, Spendezentrum and Spende. It never touches
the database directly -- it only calls the API via HTTP.

\textbf{Packaging.} The client script is wrapped into a Debian package with a
minimal \texttt{DEBIAN/control} file and installed as a system-wide command:

\begin{lstlisting}[caption={Building and installing the .deb}]
mkdir -p blutspende-client/DEBIAN
mkdir -p blutspende-client/usr/local/bin
cp client.py blutspende-client/usr/local/bin/blutspende-client
chmod +x blutspende-client/usr/local/bin/blutspende-client
dpkg-deb --build blutspende-client
sudo dpkg -i blutspende-client.deb
\end{lstlisting}

After installation, the client is available system-wide as the
\texttt{blutspende-client} command.

% ============================================================
\section{Operation and deployment}
% ============================================================

The backend runs as a PostgreSQL container orchestrated by Docker Compose. The
schema is loaded automatically from a mounted SQL file on first start.

\begin{lstlisting}[caption={docker-compose.yml (excerpt)}]
services:
  db:
    image: postgres:16
    environment:
      POSTGRES_USER: dbms10
      POSTGRES_PASSWORD: dbms10pass
      POSTGRES_DB: blutspende
    ports:
      - "5433:5432"
    volumes:
      - ./system/db/schema.sql:/docker-entrypoint-initdb.d/schema.sql
      - pgdata:/var/lib/postgresql/data

volumes:
  pgdata:
\end{lstlisting}

\textbf{Deployment steps.}
\begin{enumerate}[leftmargin=1.6em, topsep=2pt]
  \item \texttt{docker-compose up -d} -- start the database.
  \item \texttt{uvicorn main:app --host 0.0.0.0 --port 8000} -- start the API.
  \item \texttt{sudo dpkg -i blutspende-client.deb} -- install the frontend.
  \item Run \texttt{blutspende-client} to interact with the system.
\end{enumerate}

% ============================================================
\section{Reflection}
% ============================================================

The separation into database, API and CLI client made the 56-day business
rule easy to verify independently at each layer: first directly in
PostgreSQL via \texttt{psql}, then through the API with \texttt{curl}, and
finally through the client itself. The trickiest part was choosing where to
enforce the rule -- a database trigger was chosen over an application-level
check so that the constraint holds regardless of which client inserts data.
Given more time, a dedicated endpoint to check a donor's eligibility date in
advance (instead of only rejecting at insert time) would improve the user
experience.

\end{document}
